\documentclass[11pt]{article}

\usepackage[utf8]{inputenc}
\usepackage[T1]{fontenc}
\usepackage{lmodern}
\usepackage{geometry}
\usepackage{amsmath,amssymb,amsfonts,amsthm,mathtools}
\usepackage{array}
\usepackage{enumitem}
\usepackage{microtype}
\usepackage{hyperref}

\geometry{margin=1in}
\hypersetup{colorlinks=true, linkcolor=black, urlcolor=blue, citecolor=black}
\setlist[enumerate]{topsep=0.35em,itemsep=0.3em}
\setlist[itemize]{topsep=0.3em,itemsep=0.25em}

\newcommand{\EqSrc}{\mathsf{Eq}_{\mathrm{src}}}
\newcommand{\EqLC}{\mathsf{Eq}_{\mathrm{src}}^{\mathrm{lc}}}
\newcommand{\ReadLC}{\mathsf{Read}_{\mathrm{lc}}}
\newcommand{\ObsLocLC}{\mathsf{ObsLoc}_{\mathrm{lc}}}
\newcommand{\RetainH}{\mathsf{Retain}_{H}}
\newcommand{\GenH}{\mathsf{Gen}_{H}}
\newcommand{\Msrc}{\mathcal{M}_{\mathrm{src}}}
\newcommand{\RespLC}{\mathsf{Resp}_{\mathrm{lc}}}
\newcommand{\SelResp}{\mathsf{Sel}_{\RespLC}}
\newcommand{\geff}{g_{\mathrm{eff}}}
\newcommand{\Block}{\mathsf{Block}}

\newtheorem{definition}{Definition}
\newtheorem{lemma}{Lemma}
\newtheorem{theorem}{Theorem}
\newtheorem{corollary}{Corollary}
\newtheorem{payload}{New mathematical payload}
\newtheorem{distancestatus}{Distance-to-GR status}
\newtheorem{controlrule}{Control rule}

\title{\textbf{Localization Source-Basis Axiom Selector-Domain EqSrc Ledger-Complete Source-Family Generator Packet Audited ObsLoc Observer-Readout RespLC Source-Side Selector No-Go Theorem Packet}\\[0.35em]
\large Scoped Non-Determinacy of Sign, Scale, and Response-Token Selectors}
\author{Ontology Formalizer AgentJob AJ-RT-20260614-050-001}
\date{June 17, 2026}

\begin{document}

\maketitle

\begin{abstract}
This draft/control Ontology Formalizer artifact executes the packet
selected in \texttt{RT-20260614-049}. The question is whether the
non-instantiation found in \texttt{RT-20260614-048} can be upgraded
from local absence to a scoped source-side no-go theorem for the
current retained tuple
\[
    S_X=(U_X,[n]_{U_X},\Phi_{U_X},T^n,T^\Phi)
\]
and the audited no-target-import selector contract
\[
    \SelResp(S_X)=(o^R_X,\nu^R_X,\kappa^R_X).
\]
Under the explicit hypotheses below, the answer is yes in a scoped
sense: \(S_X\) alone does not determine response sign orientation,
positive response normalization, or response-token semantics. Any
selector that supplies \(o^R_X\), \(\nu^R_X\), or \(\kappa^R_X\)
must add structure outside the retained tuple or violate the
no-target-import admissibility rule. This is not a global theorem
against future source extensions. It supplies no canonical
\(\RespLC\), \(\Msrc\), \(\geff\), matter coupling, Einstein
equations, \(\EqSrc\), \(\RetainH\), \(\GenH\), finite-variation
robustness, benchmark promotion, Gate Chair review, candidate
reconstruction, completed derivation, or global theory rejection.
\end{abstract}

\section{Controlled Input}

The retained source tuple is
\[
    S_X=(U_X,[n]_{U_X},\Phi_{U_X},T^n,T^\Phi),
\]
where \(U_X=\{a_0,a_1,a_2\}\), \([n]_{U_X}\) is an exact
normal-orbit token, \(\Phi_{U_X}=\ReadLC(X_{\mathrm{loc}})|_{U_X}\),
and \(T^n,T^\Phi\) are exact-branch certificate transports.

The selector contract from \texttt{RT-20260614-046}, audited in
\texttt{RT-20260614-047}, requires source-only components
\[
    o^R_X,\qquad \nu^R_X,\qquad \kappa^R_X,
\]
for response sign orientation, positive response normalization, and
response-token semantics. The Refuter stress test in
\texttt{RT-20260614-048} found that the audited contract is clean as a
burden but constructively empty relative to the retained tuple.

\section{Hypotheses}

\begin{definition}[Admissible selector determination]
For this packet, an admissible selector determination is a rule
\[
    D(S_X)=(o^R_X,\nu^R_X,\kappa^R_X)
\]
whose inputs are limited to \(U_X,[n]_{U_X},\Phi_{U_X},T^n,T^\Phi\)
and already-tracked source-control data. The rule may use finite-set
operations, equality, relabeling-invariant comparisons, source
readout syntax, and exact-branch transport syntax. It may not use
target metric geometry, target proper time, target observer protocol,
target calibration, matter species, stress-energy response, benchmark
rank, validation status, role identity, or generated visibility.
\end{definition}

\begin{definition}[Source-invariance requirement]
An admissible determination must be invariant under transformations
that preserve the retained source tuple. If a transformation leaves
\[
    U_X,\ [n]_{U_X},\ \Phi_{U_X},\ T^n,\ T^\Phi
\]
unchanged as tracked source data, then it cannot be used by
\(D\) to distinguish response sign, scale, or token semantics.
\end{definition}

\begin{definition}[Nontrivial selector target]
A nontrivial \(\RespLC\) selector target must distinguish at least one
of the following:
\[
    \text{response sign orientation},\quad
    \text{positive response scale},\quad
    \text{response-token identity}.
\]
A rule that merely names an arbitrary convention is not a
source-determined selector unless the convention is itself sourced
from the retained tuple.
\end{definition}

\section{Source-Invisible Transformations}

\begin{lemma}[Response transformations leave the retained tuple fixed]
Let \(R=\{(\sigma_{ij},w_{ij},\rho_{ij})\}_{(i,j)\in P}\) be a
proposed response representative over source-indexed pairs. The
transformations
\[
    \sigma_{ij}\mapsto -\sigma_{ij},\qquad
    w_{ij}\mapsto \lambda w_{ij}\quad(\lambda>0),\qquad
    \rho_{ij}\mapsto \tau(\rho_{ij})
\]
for a response-token bijection \(\tau\) leave the retained source
tuple \(S_X\) fixed.
\end{lemma}

\begin{proof}
The tuple \(S_X\) contains source labels, a normal-orbit token,
restricted readout syntax, and exact-branch transports. It contains no
signed response representative, no positive response unit, and no
response-token semantics map. Changing \(\sigma_{ij}\), rescaling
\(w_{ij}\), or relabeling \(\rho_{ij}\) therefore changes only a
proposed response codomain. It does not change \(U_X\),
\([n]_{U_X}\), \(\Phi_{U_X}\), \(T^n\), or \(T^\Phi\).
\end{proof}

\begin{lemma}[The retained tuple supplies no selector-breaking datum]
No component of \(S_X\) breaks the sign-reversal, positive-rescaling,
or response-token-relabeling freedom stated in the preceding lemma.
\end{lemma}

\begin{proof}
The finite set \(U_X\) gives labels but no response orientation.
The orbit token \([n]_{U_X}\) records an equivalence class rather than
a signed representative. The readout syntax \(\Phi_{U_X}\) is not a
detector calibration or metric scale. The transports \(T^n,T^\Phi\)
preserve exact-branch certificate syntax but do not select a response
unit or matter-response token. Thus none of the retained components
distinguishes the transformed response representatives.
\end{proof}

\section{Scoped No-Go Theorem}

\begin{theorem}[Scoped source-side selector no-go theorem]
Under the admissible selector-determination and source-invariance
hypotheses above, the retained tuple
\[
    S_X=(U_X,[n]_{U_X},\Phi_{U_X},T^n,T^\Phi)
\]
does not determine a nontrivial selector
\[
    \SelResp(S_X)=(o^R_X,\nu^R_X,\kappa^R_X).
\]
Any rule that supplies response sign orientation, positive response
normalization, or response-token semantics must either add structure
not present in \(S_X\) or use a forbidden import.
\end{theorem}

\begin{proof}
Assume an admissible determination
\[
    D(S_X)=(o^R_X,\nu^R_X,\kappa^R_X)
\]
exists and is nontrivial. By nontriviality it distinguishes response
sign orientation, positive response scale, response-token identity, or
some combination of these. By the first lemma, sign reversal,
positive rescaling, and response-token relabeling leave the retained
source tuple fixed. By the second lemma, no retained source component
breaks those freedoms.

The source-invariance requirement then forces \(D\) to assign the same
source-determined value to response representatives related by those
transformations. But a nontrivial selector must distinguish at least
one of the transformed response components. This is a contradiction.

The only escape is to add a convention or datum that is not contained
in \(S_X\): for example a signed orientation convention, a response
unit, a detector calibration, a matter-response semantics map, a
target metric, an observer protocol, benchmark evidence, or an
authority label. Those escapes are either outside the retained tuple
or explicitly forbidden by the audited selector contract. Therefore
\(S_X\) does not determine \(\SelResp(S_X)\) under the stated
hypotheses.
\end{proof}

\begin{corollary}[Local absence upgraded to scoped theorem]
The non-instantiation recorded in \texttt{RT-20260614-048} is not only
an unfilled local file slot. It follows from the current retained
tuple plus the no-target-import selector contract and the invariance
requirements stated here.
\end{corollary}

\begin{corollary}[No global rejection]
The theorem is scoped to the current tuple, the audited selector
contract, and the stated admissible operations. It does not prove that
a future source-language extension, a human-gated ontology edit, or a
new source primitive could never supply selector data.
\end{corollary}

\section{Countermodel or Obstruction}

Consider two proposed response codomains \(R\) and \(R'\) related by
\[
    (\sigma_{ij},w_{ij},\rho_{ij})
    \mapsto
    (-\sigma_{ij},\lambda w_{ij},\tau(\rho_{ij})).
\]
Both codomains are compatible with the same retained tuple \(S_X\)
because \(S_X\) contains no datum distinguishing the response sign,
positive scale, or token semantics. A selector that chooses \(R\) over
\(R'\) therefore adds an extra convention. A selector that treats them
as indistinguishable does not supply \(o^R_X\), \(\nu^R_X\), or
\(\kappa^R_X\). This is the scoped obstruction.

\section{New Mathematical Payload}

\begin{payload}[Theorem with hypotheses and proof]
The theorem above proves a scoped source-side selector no-go result:
the current retained tuple cannot determine a nontrivial
\(\SelResp\) under the audited no-target-import selector contract and
source-invariance requirement.
\end{payload}

\begin{payload}[Countermodel or obstruction]
The source-invisible response-codomain transformation family supplies
a concrete obstruction: transformed response representatives disagree
on sign, scale, and token semantics while preserving all retained
source data.
\end{payload}

\section{Distance-to-GR Status}

\begin{distancestatus}[Matrix]
\begin{center}
\small
\begin{tabular}{p{0.32\linewidth}p{0.58\linewidth}}
\hline
\textbf{Burden} & \textbf{Status} \\
\hline
Source ontology primitives & Unchanged; no canonical ontology edit is made and no selector primitive is adopted. \\
Source equivalence \(\EqSrc\) & General \(\EqSrc\) remains undischarged; the theorem concerns only the current \(\RespLC\) selector burden. \\
Finite variation robustness & Not discharged; the theorem is exact-tuple scoped and does not cover relaxed boundary transport nonmatching composition omitted inverse histories quotient readout coordinates or negative-witness coverage. \\
Concrete negative witnesses & A scoped no-go obstruction is supplied for \(o^R_X,\nu^R_X,\kappa^R_X\) relative to the current \(S_X\). \\
Observer normal/readout orbit & Retained as source syntax; it does not determine response sign orientation scale normalization or token semantics. \\
Effective Lorentzian metric & Not supplied; \(\Msrc\) and \(\geff\) remain blocked. \\
Universal matter coupling & Not supplied; response-token semantics remain unavailable from current source data. \\
Einstein equations & Not supplied. \\
Benchmark promotion & Blocked. \\
Gate Chair review & Blocked and human-gated. \\
Current line hard-fail & Scoped no-go theorem supplied for the current selector path only; no global theory rejection or future source-extension impossibility is claimed. \\
\hline
\end{tabular}
\end{center}
\end{distancestatus}

\section{Parent-Child Synthesis}

This artifact is the fused parent output for
\texttt{role\_decomposition.mode = parent\_child\_parallel\_synthesis}.
The child outputs are:
\begin{itemize}
    \item \texttt{child\_phys\_math\_resp\_lc\_source\_side\_selector\_no\_go\_theorem.yaml};
    \item \texttt{child\_phys\_phil\_resp\_lc\_source\_side\_selector\_no\_go\_theorem.yaml}.
\end{itemize}
The conflict review is
\texttt{parent\_conflict\_review\_resp\_lc\_source\_side\_selector\_no\_go\_theorem.yaml}
and found no blocking conflict. The fusion notes are
\texttt{parent\_fusion\_notes\_resp\_lc\_source\_side\_selector\_no\_go\_theorem.md}.

\section{Control Boundary}

\begin{controlrule}[Scoped result only]
This theorem applies only to the current retained tuple, the audited
no-target-import selector contract, and the admissible operations
stated in this packet. It may not be cited as a global impossibility
result for all future source extensions.
\end{controlrule}

\begin{controlrule}[No promotion]
The theorem supplies no \(\RespLC\) adoption, no \(\Msrc\), no
\(\geff\), no matter coupling, no Einstein equations, no benchmark or
Gate Chair authority, no candidate reconstruction, and no
completed-derivation language.
\end{controlrule}

\begin{controlrule}[Next-route discipline]
The logical next step is a bounded Refuter stress test of the scoped
no-go theorem. The stress test should examine hypothesis strength,
arbitrary-convention loopholes, possible source-extension
countermodels, and overread risk.
\end{controlrule}

\section*{APA 7 References}

Ricciardi, A. S. (2026, June 17). \emph{RespLC theoretical
continuation selector decision} [Unpublished control artifact]. The
Aether-Flow Interpretation of Relativity Research Project,
\texttt{research\_control/tasks/RT-20260614-049/artifacts/90\_RESP\_LC\_THEORETICAL\_CONTINUATION\_SELECTOR\_DECISION.yaml}.

Ricciardi, A. S. (2026, June 17). \emph{Localization source-basis
axiom selector-domain EqSrc ledger-complete source-family generator
packet audited ObsLoc observer-readout RespLC source-side selector
instantiation Refuter stress test: Scoped obstruction for
no-target-import selector instantiation} [Unpublished manuscript].
The Aether-Flow Interpretation of Relativity Research Project,
\texttt{research\_control/tasks/RT-20260614-048/artifacts/89\_LOCALIZATION\_SOURCE\_BASIS\_AXIOM\_SELECTOR\_DOMAIN\_EQSRC\_LEDGER\_COMPLETE\_SOURCE\_FAMILY\_GENERATOR\_PACKET\_AUDITED\_OBSLOC\_OBSERVER\_READOUT\_RESP\_LC\_SOURCE\_SIDE\_SELECTOR\_INSTANTIATION\_REFUTER\_STRESS\_TEST.tex}.

Ricciardi, A. S. (2026, June 17). \emph{Localization source-basis
axiom selector-domain EqSrc ledger-complete source-family generator
packet audited ObsLoc observer-readout RespLC source-side selector
primitive smuggling audit: Hidden-import screens for the draft
selector contract} [Unpublished manuscript]. The Aether-Flow
Interpretation of Relativity Research Project,
\texttt{research\_control/tasks/RT-20260614-047/artifacts/88\_LOCALIZATION\_SOURCE\_BASIS\_AXIOM\_SELECTOR\_DOMAIN\_EQSRC\_LEDGER\_COMPLETE\_SOURCE\_FAMILY\_GENERATOR\_PACKET\_AUDITED\_OBSLOC\_OBSERVER\_READOUT\_RESP\_LC\_SOURCE\_SIDE\_SELECTOR\_PRIMITIVE\_SMUGGLING\_AUDIT.tex}.

Ricciardi, A. S. (2026, June 17). \emph{Localization source-basis
axiom selector-domain EqSrc ledger-complete source-family generator
packet audited ObsLoc observer-readout RespLC source-side selector
primitive packet: Draft/control contract for sign orientation, scale
normalization, and response-token semantics} [Unpublished manuscript].
The Aether-Flow Interpretation of Relativity Research Project,
\texttt{research\_control/tasks/RT-20260614-046/artifacts/87\_LOCALIZATION\_SOURCE\_BASIS\_AXIOM\_SELECTOR\_DOMAIN\_EQSRC\_LEDGER\_COMPLETE\_SOURCE\_FAMILY\_GENERATOR\_PACKET\_AUDITED\_OBSLOC\_OBSERVER\_READOUT\_RESP\_LC\_SOURCE\_SIDE\_SELECTOR\_PRIMITIVE\_PACKET.tex}.

\end{document}
